% Research note: closing helices at quantization time.
% Build: pdflatex paper && bibtex paper && pdflatex paper && pdflatex paper
\documentclass[11pt,a4paper]{article}

\usepackage[T1]{fontenc}
\usepackage[utf8]{inputenc}
\usepackage{lmodern}
\usepackage[margin=2.5cm]{geometry}
\usepackage{amsmath,amssymb,amsthm}
\usepackage{booktabs}
\usepackage[numbers,sort&compress]{natbib}
\usepackage[hidelinks]{hyperref}

\newtheorem{definition}{Definition}
\newtheorem{proposition}{Proposition}
\newtheorem{lemma}{Lemma}
\newtheorem{corollary}{Corollary}
\newtheorem{remark}{Remark}
\newtheorem{conjecture}{Conjecture}
\newtheorem{problem}{Open problem}

\newcommand{\Z}{\mathbb{Z}}
\newcommand{\R}{\mathbb{R}}
\newcommand{\shift}{s}

\title{Closing Edge Loops at Quantization Time:\\
Helix Constraints for T-Mesh Quantization\\[0.5em]
\large Research note --- draft, not peer reviewed}
\author{openclare project}
\date{17 October 2026}

\begin{document}
\maketitle

\begin{abstract}
Automatic quad remeshers produce edge loops that wind round a part of the
shape as helices instead of closing into rings, a defect that makes such
meshes awkward to edit by hand. Bommes et
al.~\cite{bommes2011structure} identified these \emph{q-helices} and remove
them after meshing, by grid-preserving operators applied to the finished
quad mesh. We take the opposite route and prevent them where the integer
structure of the mesh is decided: in the quantization of a T-mesh. A
straight walk across the rectangles of a T-mesh that comes back to its
start defines a \emph{turn}; we show that the pitch of the corresponding
helix is a linear form in the integer lengths of the arcs, so that closing
the helix is a linear constraint on the quantization. We add such
constraints to a minimum-deviation-flow quantizer~\cite{heistermann2023bimdf}
through a least-change projection and targeted repair walks, select among a
few ways of finding turns by the share of spiral edges they leave, and
guard against folds and budget drift. On our bench the share of edges on
spiral loops at 3{,}000 quads drops from 16.7\% to 6.0\% on twelve models;
on a human body remeshed at the face count of the hand-made base mesh it
drops from 12.1\% to 1.1\%, below the 2.5\% of the hand-made mesh. We relate
turns to the twist of flat cylinders in the theory of Jenkins--Strebel
differentials~\cite{strebel1984quadratic}, show with the holonomy of the
flat metric that a pair of singularities of valences three and five lying on
one trace between two turns forbids closing both, and list the open problems
that a complete treatment would need to solve.
\end{abstract}

\section{Introduction}

Field-aligned quad remeshing has reached a level where element shape and
alignment are rarely the limiting factor. What remains is structure: the
way edge loops run across the surface. Artists expect a loop round an arm,
a leg or a neck to be a ring. Automatic methods often produce a helix
instead: the loop goes round, comes back one or a few rows further, and
winds along the whole part. The problem is old and well identified:
Bommes et al.\ showed in 2011 that helical configurations are what keeps
the base complex of an otherwise good quad mesh dense~\cite{bommes2011structure}.

This note documents an approach developed in the open-source remesher
openclare, and positions it against the literature we could check. Our
contributions are:
\begin{enumerate}
\item \emph{Turns} of a T-mesh and their \emph{shift}: a straight walk
  across the rectangles of a T-mesh that returns to its starting side,
  and a linear form in the arc lengths whose integer value is the pitch of
  the helix the walk becomes in the quad mesh (Section~\ref{sec:turns}).
\item Closing turns within T-mesh quantization: a least-change projection
  of the target lengths, an integer search, and repair walks judged on the
  turns, with safeguards on folds, distortion and budget
  (Section~\ref{sec:closing}).
\item A classification of the edge loops of a quad mesh into closed,
  spiral, open and border loops, with spirals of order up to four
  (Section~\ref{sec:measure}), and a bench that holds automatic remeshers
  against meshes modelled by hand (Section~\ref{sec:experiments}).
\item Obstructions to closing: the shift of a turn is the Burgers vector of
  the singularities on one side; a pair of valences three and five joined by
  a trace between two turns forces one of them to be a helix; and a
  checkerboard rule gives the parity of the shift between two turns
  (Section~\ref{sec:obstructions}).
\end{enumerate}
The results are preliminary: one seed per model, one implementation, and
heuristics where an exact formulation is still missing
(Section~\ref{sec:open}).

\section{Background and related work}\label{sec:related}

\paragraph{Counting singularities.}
On a closed surface of genus $g$, the valences $k(v)$ of a quad mesh
satisfy $\sum_v (4 - k(v)) = 4\chi = 8(1-g)$. With only valences 3 and 5,
$n_3 - n_5 = 8(1-g)$: a sphere or a closed human body needs at least eight
vertices of valence three. This bound says how many singular vertices are
needed, not where they go; a hand-made human base mesh has hundreds.

\paragraph{Quad meshes as quartic differentials.}
Lei et al.~\cite{lei2020quartic} show that a closed quad mesh induces a
Riemann surface and a meromorphic quartic differential $\omega = (dz)^4$
whose zeros and poles are the singular vertices, a vertex of valence $k$
being of order $k-4$; conversely a quartic differential with finite
trajectories induces a quad mesh. The divisor of singularities then
satisfies the Abel--Jacobi condition, and a divisor satisfying it is the
divisor of a meromorphic quartic differential. Zheng et
al.~\cite{zheng2020abeljacobi} optimize the singularity configuration by
integer programming to satisfy the condition. Two points matter here.
First, the condition is stated with holomorphic one-forms, which a surface
of genus zero does not have; our reading is that in genus zero it adds
nothing to the degree condition (this is an inference, not a statement of
the papers). Second, the paper lists as an open problem the guarantee that
all trajectories are \emph{finite}. A trajectory that is not finite is
precisely an edge loop that does not close: the helices of this note.

\paragraph{Closed trajectories.}
A quadratic differential all of whose non-critical horizontal trajectories
are closed is a Jenkins--Strebel differential~\cite{strebel1984quadratic};
the complement of its critical graph is a union of flat cylinders, each
with a height and a \emph{twist}, the shift with which the cylinder is glued
back after a cut along its core curve. In a quad mesh, a ring of quads is
such a cylinder; a helix is a cylinder whose gluing is shifted by a whole
number of rows. Section~\ref{sec:turns} makes this precise at the level of a
T-mesh.

\paragraph{Quad meshes as flat metrics.}
Declaring every quad a unit square turns a quad mesh into a flat metric
with cone points at its irregular vertices. Chen et al.~\cite{chen2019metric}
characterize the metrics that arise this way by four conditions:
Gauss--Bonnet, rotational holonomy in multiples of a quarter-turn,
alignment with the border, and a \emph{finite geodesic lamination}: the
straight lines of the field are finite loops. On a surface of genus zero
the holonomy condition follows from Gauss--Bonnet (their Lemma~3.9). The
translational part of the holonomy is not discussed, and the finiteness
condition is assumed, not constructed. Lei et
al.~\cite{lei2023crossfields} show that the singularity configurations of
quad meshes are a strict subset of those of cross fields, and explain in
this way why no quad mesh of the torus has exactly one vertex of valence
three and one of valence five.

\paragraph{Dislocations.}
Izmestiev et al.~\cite{izmestiev2013torus} give geometric proofs of such
non-existence results, among them the torus with one vertex of valence three
and one of valence five, first proved combinatorially by Barnette, Jucovi\v{c}
and Trenkler in 1971 (as cited there). Their tool is the holonomy of the
developing map of the flat metric, a rigid motion per loop. When the
rotation of a loop is the identity, its translation is the \emph{Burgers
vector} of crystallography; two opposite cone points at distinct positions
always have a non-zero one, which is the dislocation a pair of valences three
and five makes in the square grid. Section~\ref{sec:obstructions} uses this
tool for the shifts of turns.

\paragraph{Helices in quad meshes.}
Bommes et al.~\cite{bommes2011structure} define a \emph{q-helix} as a dual
parametric line of a quad mesh with an integer \emph{pitch} $h$ and
\emph{turn length} $\tau$: walking $\tau$ steps along it is equivalent to
$h$ side-steps. They detect minimal q-helices ($h \le \tau$, no sub-helix of
smaller pitch) by the self-intersections of orthogonal dual lines, and
remove them greedily with \emph{grid-preserving operators}: closed dual
paths of shift and collapse steps, found by a graph search, that keep the
mesh all-quad and its singularities in place. Tangential smoothing then
repairs the distortion the shifts introduce. Their experiments remove most
helices (from 19--52 down to 1--23 per model) and coarsen the base complex
accordingly. They note that aligning pairs of singularities directly in the
parametrization, as Myles et al.\ do for T-meshes~\cite{myles2010feature},
is ambiguous to automate. The approach works on a finished mesh; it does not
touch the parametrization or its quantization.

\paragraph{T-mesh quantization.}
Integer-grid maps~\cite{bommes2009miq} and their quantization decide the
integer structure of a quad mesh. Campen et al.\ quantize via the
integer lengths of the arcs of a T-mesh traced in a seamless
parametrization~\cite{campen2015qgp}. Lyon et al.\ solve a constrained
version as an integer linear program, bounding the directional deviation of
the separatrices of the resulting layout~\cite{lyon2021constrained} (we
could only read the abstract). Heistermann et al.\ cast quantization as a
minimum-deviation flow on a bi-directed graph and give a fast approximate
solver and an exact refinement~\cite{heistermann2023bimdf}; openclare
follows their formulation. To our knowledge none of these constrains the
helices: the lengths are chosen for consistency, validity and closeness to
the parametrization, and the loops close or spiral as a side effect.

\paragraph{Periods round a tube.}
QuadCover~\cite{kalberer2007quadcover} lifts the cross field to a branched
covering where it becomes a vector field, and integrates it there. The
stripe parameterization of tubular surfaces~\cite{kalberer2011stripe}, built
on it, measures the jump (``gap'') of the parameter across each cut path,
which is constant along the path, and rounds it to the nearest whole number;
the grid lines are continuous exactly when every gap is whole. The gap of
the coordinate across a loop round a tube is what we call the shift of a
turn: rounded to zero it gives rings, rounded to a non-zero whole number a
helix of that pitch. That paper treats stripes spiralling round the tube as
a legitimate outcome, and notes that rounding the gaps for least distortion
is a hard combinatorial problem that it does not solve exactly.

\paragraph{Layouts traced from the field.}
Myles et al.~\cite{myles2014robust} trace the field into a partition into
quadrilateral patches and adjust the layout until its lengths are
consistent; Pietroni et al.~\cite{pietroni2016tracing} select traced
separatrices by a binary linear program; Razafindrazaka et
al.~\cite{razafindrazaka2015perfect,razafindrazaka2017optimal} obtain a
layout as a minimum-weight perfect matching on a graph of the singularities.
Campen surveys these methods~\cite{campen2017survey}. They build coarse
layouts; whether the dense mesh made from the layout spirals is, as far as
the abstracts we read state, not among their criteria.

\paragraph{Lower bounds.}
We found no result that bounds from below the number of helices of a quad
mesh, or the number of its loops that do not close, for a given singularity
configuration. The nearest results are existence and non-existence theorems:
Strebel's existence of Jenkins--Strebel differentials with prescribed
cylinders on a fixed Riemann surface~\cite{strebel1984quadratic}, the
non-existence results above~\cite{izmestiev2013torus,lei2023crossfields},
and the finiteness condition of~\cite{chen2019metric}, which the open
problem of~\cite{lei2020quartic} asks to guarantee.

\paragraph{Loops by construction.}
Dual loops meshing~\cite{campen2012dual} builds a layout from loops traced
on the surface, which close by construction; it produces coarse layouts
rather than dense field-aligned meshes.

\paragraph{Measuring loops.}
QuadGPT~\cite{liu2025quadgpt} counts \emph{quad rings} (closed loops of
faces) and \emph{quad lines} (open strips) and uses the share of faces in
rings and the mean length of lines as reward signals for preference
optimization; it does not report them as evaluation metrics, nor does it
tell spirals from open lines.

\section{Turns of a T-mesh}\label{sec:turns}

\paragraph{Setting.}
A T-mesh is a partition of the surface into rectangles, the
\emph{patches}, whose sides are made of \emph{arcs}; nodes where a trace
stops on another are T-junctions, and the singular vertices of the field
are nodes. Each patch $P$ has local coordinates with its first corner at
the origin and its sides along the axes. A quantization assigns every arc
$a$ a length $\ell_a \in \Z_{\ge 0}$ such that opposite sides of every patch
have equal length; the quad mesh is then read off as the integer grid of
each patch.

\begin{definition}[Turn]
A \emph{turn} $\gamma$ is a straight walk across patches: it enters a patch
through a side, leaves through the opposite side at the same distance from
the corresponding corner, enters the neighbouring patch across the arc it
reaches, and so on, until it comes back into its first patch through its
first side, with the four quarter-turns of the crossings summing to zero.
\end{definition}

Crossing from patch $P$ to patch $Q$ across an arc composes a quarter-turn
rotation of the local frame with a translation whose components are signed
sums of arc lengths along the rims of $P$ and $Q$: the position of a rim
point is the sum of the lengths of the arcs between it and the first corner.

\begin{proposition}[The shift of a turn is linear]\label{prop:linear}
For a turn $\gamma$ there are integer coefficients $c_a$, independent of
the lengths, such that after a full turn a point of the first side has
moved across the walking direction by
\[
  \shift_\gamma(\ell) = \sum_a c_a\, \ell_a .
\]
For integer lengths, the dual line of the quad mesh that follows $\gamma$
comes back $\shift_\gamma(\ell)$ rows beside itself: it closes into a ring
if and only if $\shift_\gamma(\ell) = 0$, and is otherwise a helix of pitch
$|\shift_\gamma(\ell)|$ in the sense of~\cite{bommes2011structure}.
\end{proposition}

\begin{proof}[Sketch]
Each crossing is a map $x \mapsto R_k x + t$ with $R_k$ a rotation by $k$
quarter-turns and $t$ a vector whose entries are sums of $\pm\ell_a$. The
composition along $\gamma$ has total rotation zero, hence is a translation
$T_\gamma(\ell)$, linear in $\ell$ with integer coefficients. Its component
along the walking direction does not move a point of the first side off
that side; its component across is $\shift_\gamma(\ell)$. The walk crosses
no singularity (singularities are nodes, never inside a patch), so the
quad mesh is a regular grid along it, and a dual line at distance $o$ from
the first corner comes back at distance $o + \shift_\gamma(\ell)$.
\end{proof}

In openclare (\texttt{quantization/tubes.rs}) the shift is built once per
turn as an expression in the arc lengths, by composing the rim expressions
of the crossed patches; evaluating it with the real lengths of the
parametrization gives the \emph{real shift}, the twist the field itself
would give the cylinder.

\begin{definition}[Strip]
For whole lengths $\ell$ balancing every patch, the quad mesh is, in each
patch, the integer grid of its sides. A \emph{strip} is a maximal sequence
of quads, each sharing with the next the edge opposite the one it shares
with the previous.
\end{definition}

\begin{proposition}[Completeness at given lengths]\label{prop:complete}
Let $\ell$ be whole lengths balancing every patch. Every strip of the quad
mesh that crosses a side $\sigma$ of a patch $P$ at the quad edge
$[k, k+1]$ of $\sigma$ is, as long as it stays in the regular region, the
walk of Definition~1 taken \emph{with the lengths $\ell$} from $\sigma$ at
offset $k + \tfrac12$. In particular, walking from every half-integer of
every side finds every strip of the quad mesh that comes back to the side
it left; by Proposition~\ref{prop:linear} it is a helix exactly when
$\shift_\gamma(\ell) \neq 0$.
\end{proposition}

\begin{proof}
Inside $P$, the quads of the strip are the unit squares between the grid
lines at $k$ and $k+1$ across $\sigma$; their dual line is the straight
line at $k + \tfrac12$, which reaches the opposite side at the same distance
from the corresponding corner, which is the move of the walk. At the arc
where it leaves $P$, the crossing map to the neighbouring patch is a
quarter-turn followed by a translation whose components are sums of
$\pm\ell_a$, hence integers: it maps the grid of $P$ onto the grid of the
neighbour, and the midpoint of a quad edge onto the midpoint of a quad
edge. Every node lies at a whole distance along the sides, so a walk at a
half-integer offset never meets a node, and the arc it crosses is
determined. The strip and the walk thus continue together, patch after
patch. A
strip only meets singular vertices at its sides, never inside, since they
are corners of patches; the regular region is where the grid coordinates
of the patches glue by such maps.
\end{proof}

The sampled walks of Section~\ref{sec:closing} start at fixed shares of
the real side lengths. Proposition~\ref{prop:complete} says where the
helices of a given mesh are \emph{all} found: at the half-integers of the
whole lengths. This suggests closing by \emph{cuts}, in the sense of
cutting planes:

\begin{enumerate}
\item quantize; let $\ell$ be the whole lengths;
\item walk every strip at $\ell$ (Proposition~\ref{prop:complete}); among
  the ways round that come back, select those meant to be rings;
\item add their shift expressions to the closing constraints, keeping those
  already closed, as long as the bounds of separation allow it;
\item solve again in whole numbers; repeat until no selected way is a
  helix.
\end{enumerate}

\begin{remark}
A turn round a part of the shape whose real shift is far from zero is not a
rounding artefact: the field winds round the part. We only hold turns whose
real shift is a small share of their length, the \emph{lean}; a turn round
a group of singularities, by contrast, may legitimately come back shifted.
\end{remark}

\section{Closing turns in quantization}\label{sec:closing}

The quantizer of openclare minimizes, as in~\cite{heistermann2023bimdf}, a
weighted deviation $\sum_a w_a\,|\ell_a - \lambda\, r_a|$ of the integer
lengths from the scaled real lengths $\lambda r_a$, subject to the balance of
each patch, by cancelling improving walks in a strip network; the scale
$\lambda$ is searched for the number of quads asked for, and arcs that would
merge two singularities are held open. Closing turns is a side constraint
that is not a flow constraint. We handle it in three steps.

\paragraph{Finding turns.}
From points at given shares along each side of each patch, we walk straight
across patches and keep the walks that return. A turn is \emph{held} if its
real shift is above a small threshold and its lean below a bound, and if
its length is below a multiple of the side of the T-mesh (the side of a
square of its area). Several \emph{ways} are tried: starting shares
$\{0.5\}$ or $\{0.2, 0.5, 0.8\}$, lean bounds $0.05$ or $0.12$, length
bounds $2$ or $6$ sides.

\paragraph{Least-change projection.}
The target lengths are changed as little as possible, in the weighted least
squares sense, for every held turn to have zero shift and every patch to
stay balanced:
\[
  \min_{x}\ \sum_a \frac{(x_a - r_a)^2}{w_a}
  \quad\text{s.t.}\quad A x = 0,\quad \shift_\gamma(x) = 0\ \ \forall\gamma\in H,
\]
solved by a sparse Cholesky factorization of $A W^{-1} A^\top$. If the
system cannot be met, the held turn left with the largest shift is let go
and the projection repeated.

\paragraph{Integer search and repair.}
The integer search of the quantizer runs on the projected targets. Rounding
leaves some turns shifted by one or a few rows. Repair walks are then taken
one at a time: a walk is guided by a bias along the arcs of the shifted
turns, and accepted if it reduces the total shift, costs little in
deviation, keeps the structure intact and the number of quads within
tolerance. A second, \emph{aimed} repair treats each held turn still shifted
on its own: the bias is put on its arcs alone, and a walk is accepted if it
shifts that turn less and the others no more.

\paragraph{Selection and safeguards.}
Each way gives a quantization. Its quad mesh connectivity is built and the
share of edges on spiral loops measured (Section~\ref{sec:measure}); the
way with the fewest spirals is kept, provided its number of quads is within
tolerance and its mean deviation from the real lengths at most $1.5$ times
that of the plain search. The ways are independent and run in parallel.
After meshing, if the kept quantization makes a face fold over the surface,
the plain quantization is meshed as well and the mesh with fewer folds kept.

\section{Measuring loops}\label{sec:measure}

An edge loop is followed straight through every regular vertex (valence
four, interior), each edge belonging to one loop. A loop is
\emph{closed} if it comes back to its start; \emph{border} if it runs from
border to border; \emph{open} if it stops at an irregular vertex or a
non-quad face; and \emph{spiral}, whatever its ends, if a strip of quads
leaving it across meets it again, the same way round, within
$k_{\max} = 4$ quads. The last test catches helices of pitch up to four,
including bundles of interleaved helices~\cite{bommes2011structure}, while a
ring is met again only round a whole torus or over a cap the other way
round. We report the share of edges on each kind of loop, edges of borders
left out.

\section{Experiments}\label{sec:experiments}

All numbers come from the openclare bench (repository
\texttt{marcgardent/openclare}), one seed per model, measured on the
output meshes against their sources. Models: twelve organic and mechanical
models for development, a corpus of 31 models for the full bench, and two
human bodies from Blender's Essentials library (CC0) whose base meshes,
modelled by hand, serve as references.

\paragraph{Spirals on the twelve development models, 3{,}000 quads.}
Table~\ref{tab:dev} shows the successive steps. Fidelity (Hausdorff
distance, per mil of the diagonal) and the mean corner deviation stay
within about one per cent; the last step trades a little of both for a
third fewer spiral edges.

\begin{table}[h]
\centering
\begin{tabular}{lrrr}
\toprule
& spiral edges & Hausdorff (\textperthousand) & corner deviation (deg) \\
\midrule
no closing & 16.7\% & 8.77 & 4.75 \\
closing, one way kept by spirals & 9.7\% & 8.57 & 4.85 \\
+ leaning ways, aimed repair & 9.1\% & 8.58 & 4.91 \\
+ longer turns, deviation cap & 6.0\% & 8.68 & 5.01 \\
\bottomrule
\end{tabular}
\caption{Mean over twelve models; each line adds to the one above.}
\label{tab:dev}
\end{table}

\paragraph{Against other remeshers.}
On the full bench at 3{,}000 quads, the mean share of spiral edges is
19.1\% for Quad Remesher, 21.9\% for QuadriFlow and 25.3\% for AutoRemesher,
against 12.3\% for openclare with the first version of closing (and 17.3\%
without).

\paragraph{Against the hand-made mesh.}
The male body of the Essentials library has a base mesh of 10{,}590 quads
modelled by hand, with 2.5\% of spiral edges. Remeshed by openclare at that
count before closing, it had 12.1\% (one leg was a single helix from hip to
ankle, its turns leaning more than the first lean bound allowed). With the
method of Section~\ref{sec:closing}: 10{,}467 quads, 1.1\% spiral edges,
36.2\% of loops closed against 24.7\% for the hand-made mesh, a Hausdorff
distance of 4.6\textperthousand\ against 6.5\textperthousand, and a mean
corner deviation of 4.1 degrees against 9.2.

\paragraph{The quad sphere.}
Not a result of closing but of the field. An icosphere remeshed at
3{,}000 quads had 32 irregular vertices (20 of valence three, 12 of valence
five), the pattern of its own facets, because the field was pulled towards
curvature directions that a sphere does not have; at 1{,}000 quads it had
the eight poles but no structure (a base complex of 1{,}010 patches) and
87\% of spiral edges. With directions not trusted where the two principal
curvatures are within 10\% of one another, both give the layout of a cube:
eight poles of valence three, six patches, no spiral.

\section{An exact look at the closing system}\label{sec:exact}

To see how far the projection of Section~\ref{sec:closing} is from what
the constraints allow, openclare exports, for each way of finding turns,
the system it works on: the real lengths of the 2{,}069 arcs of the T-mesh
of the male body at 10{,}590 quads, the 1{,}552 balance equations of its
patches, every turn found with its shift, and the lower bounds the integer
search placed on arcs to keep singularities apart (30 to 36 arcs held at
length one). We solve it with HiGHS through SciPy, in real numbers.

\paragraph{Where the conflict lies.}
Every held turn can be closed on its own with all bounds kept. All held
turns can be closed together if the lengths are only asked to be
non-negative. With the bounds of separation, they cannot. Giving up the
least total amount of those bounds needed to close all held turns, a
linear program, shows that for three of the four ways examined a
\emph{single} arc must collapse, in two of them the same arc, $0.08$ quad
edges long, at the left wrist; the fourth needs eleven, round the eyes and
the hands. The conflict between closing loops and keeping singularities apart
is thus real, but very local.

\paragraph{How far the heuristic is.}
With every bound kept, the largest set of held turns that can be closed
together is found by a mixed-integer program with one binary per turn
(Table~\ref{tab:exact}). The projection keeps far fewer: it removes, one at
a time, the turn left most shifted by a least-squares projection whose
negative lengths are clamped to zero, and a single clamped arc can make it
drop dozens of turns that conflict with nothing.

\begin{table}[h]
\centering
\begin{tabular}{lrrr}
\toprule
way & held & kept by the projection & closable together, bounds kept \\
\midrule
7 & 169 & 49 & 160 \\
4 & 160 & 48 & 154 \\
5 & 78 & 46 & 73 \\
1 & 88 & 53 & 87 \\
0 & 84 & 49 & 83 \\
2 & 86 & 47 & 82 \\
3 & 52 & 46 & 51 \\
6 & 48 & 42 & 47 \\
\bottomrule
\end{tabular}
\caption{Male body at 10{,}590 quads. Real numbers; the numbering of the
ways is the order in which they finished.}
\label{tab:exact}
\end{table}

In whole numbers, the same mixed-integer program did not finish within
fifteen minutes on the two systems tried.

\paragraph{From reals to whole numbers.}
We then fed the engine better targets and measured what survives its
integer search (Table~\ref{tab:rounding}). Aiming at the lengths of the
linear program closes fewer loops in the end: the weighted $L^1$ objective
puts the changes on few arcs, the deviation from the real lengths grows,
and the variants with the fewest spirals are refused by the deviation cap.
A projection that holds lengths at a floor instead of clamping them (an
active-set method on the same Cholesky solver) keeps every held turn closed
in real numbers; yet after the integer search 117 of 169 are shifted again,
and the mesh spirals more. With the current projection, too, about 40\% of
the turns it kept are shifted after the integer search. The bottleneck is
not the projection but the rounding: the integer search follows the targets
but ignores the shifts.

\begin{table}[h]
\centering
\begin{tabular}{lrrr}
\toprule
targets of the integer search & turns closed in reals & still shifted after & spiral edges (mesh) \\
\midrule
projection, clamped (current) & about 50 & 10--23 of them & 1.1\% \\
linear program, $L^1$ & up to 160 & --- & 11.9\% \\
active set, floor 0 & all held (up to 169) & 117 of 169 & 3.3\% \\
active set, floor half the real length & 36--72 & --- & 9.8\% \\
\bottomrule
\end{tabular}
\caption{Male body at 10{,}590 quads, one run each; the mesh keeps the way
with the fewest spirals, so its share is that of a single way.}
\label{tab:rounding}
\end{table}

\section{When rings cannot all close}\label{sec:obstructions}

Section~\ref{sec:exact} found that closing all held turns conflicts with
keeping singularities apart, and that the conflict comes down to one or a
few arcs. This section gives a mechanism for such conflicts, with proofs. The
tool is the holonomy of the flat metric, as used by Izmestiev et
al.~\cite{izmestiev2013torus}; its ingredients are classical, and we know of
no earlier statement of the propositions below in terms of helices.

\paragraph{Flat metric and holonomy.}
Lengths $\ell > 0$ balancing every patch make each patch a Euclidean
rectangle. Glued along their arcs, whose lengths agree on both sides, the
rectangles give the surface $\Sigma$ a flat metric with cone points at the
singular vertices: a vertex of valence $k$ has angle $k\pi/2$ and curvature
$(4-k)\pi/2$; a T-junction has angle $2\pi$ and is not a cone point. Let
$\Sigma^\circ$ be $\Sigma$ without its singular vertices, $x_0$ a point inside
a patch, and $R$ the quarter-turn. Developing a loop $\gamma$ of $\Sigma^\circ$
based at $x_0$ patch after patch, as in the proof of
Proposition~\ref{prop:linear}, gives a motion $\hat h(\gamma)\colon x \mapsto
R^j x + t$ of the coordinates of the first patch, its \emph{holonomy}. It
depends only on the homotopy class of $\gamma$ in $\Sigma^\circ$, and $\hat h$
is a homomorphism of the fundamental group into the group of such
motions~\cite{izmestiev2013torus}. A small loop round a singular vertex $c$
of valence $k$, reached along a path that develops it at the point $p$, has
holonomy $x \mapsto R^{4-k}(x - p) + p$. When the rotation of $\hat h(\gamma)$
is the identity, its translation is the \emph{Burgers vector} of $\gamma$.

For a turn $\gamma$, let $\bar\gamma$ be the closed curve that follows the
walk and, back in its first patch, returns to its starting point along a
segment parallel to the first side, just inside the patch. Let $e_1$ be the
walking direction and $e_2$ the direction along the first side. By the proof
of Proposition~\ref{prop:linear}, $\hat h(\bar\gamma)$ is a translation and
\[
  \shift_\gamma(\ell) = \langle e_2,\ \hat h(\bar\gamma) \rangle .
\]
Shifts, developed positions and Burgers vectors are linear forms in $\ell$
with integer coefficients, as long as the walks and paths cross the same
patches. An identity between them that holds on an open set of balanced
positive lengths therefore holds for all balanced lengths, and for whole
lengths the developed positions of the nodes are points of $\Z^2$.

\begin{lemma}[Gauss--Bonnet for turns]\label{lem:gauss}
Let $\gamma$ be a turn whose curve $\bar\gamma$ is simple. If $\bar\gamma$
bounds a disc $D \subset \Sigma$, the singular vertices inside $D$ satisfy
$\sum_{c \in D} (4 - k(c)) = 4$. If the curves of two turns are disjoint and
bound an annulus $A$, then $\sum_{c \in A} (4 - k(c)) = 0$.
\end{lemma}

\begin{proof}
For a region of a flat surface with cone points, bounded by straight
segments, Gauss--Bonnet reads $\sum_c \kappa_c + \sum_i \theta_i =
2\pi\chi$, where $\kappa_c$ are the curvatures of the cone points inside and
$\theta_i$ the turning angles at the corners of the border. The curve
$\bar\gamma$ is straight except at the two ends of its closing segment, where
it turns by $+\pi/2$ and $-\pi/2$ (or not at all if the shift is zero); the
turning angles sum to zero on each curve, seen from either side. With
$\kappa_c = (4 - k(c))\pi/2$, $\chi(D) = 1$ and $\chi(A) = 0$ the two
statements follow.
\end{proof}

On a surface of genus zero, such as a body, every simple closed curve
bounds a disc on each side. Every turn with a simple curve thus splits the
singular vertices into two groups of total index four each, consistently
with the total $8(1-g) = 8$ of Section~\ref{sec:related}: a fingertip beyond
a ring holds the equivalent of four vertices of valence three. A loop round
a group of singularities whose indices do not add up to four is never a
turn, whatever the lengths. A side holding at most one singular vertex has
total index $0$ or $4 - k \neq 4$, so each side of a turn holds at least two.

\begin{proposition}[The shift is set by the singularities]\label{prop:burgers}
Let $\gamma$, $\gamma_1$, $\gamma_2$ be turns with simple curves.
\begin{enumerate}
\item[(a)] If $\bar\gamma$ bounds a disc $D$ holding the singular vertices
  $c_1, \dots, c_m$, then $\shift_\gamma(\ell) = \langle e_2, B_D \rangle$
  where $B_D$, the Burgers vector of $\bar\gamma$, is the translation of the
  product of the holonomies of small loops round $c_1, \dots, c_m$, taken in
  a suitable order along paths inside $D$:
  \[
    B_D = \sum_{i=1}^m Q_i\,(I - R_i)\,p_i ,
    \qquad R_i = R^{4 - k(c_i)},
  \]
  with $p_i$ the developed position of $c_i$ and $Q_i$ the product of the
  rotations $R_j$ of the vertices that come after $c_i$ in that order.
\item[(b)] Let $\bar\gamma_1$ and $\bar\gamma_2$ be disjoint and bound an
  annulus $A$, round which the two turns go the same way, and let a path
  $\kappa$ inside $A$ from $\bar\gamma_1$ to $\bar\gamma_2$ carry the walking
  direction of $\gamma_2$ onto that of $\gamma_1$ (the turns are
  \emph{parallel}); measure both shifts along the direction $e_2$ carried
  by $\kappa$. Then
  \(
    \shift_{\gamma_1}(\ell) - \shift_{\gamma_2}(\ell) = \langle e_2, B_A \rangle,
  \)
  where $B_A$ is the Burgers vector of the border of the disc $A$ cut along
  $\kappa$, given by the same sum over the singular vertices inside $A$,
  developed inside the cut disc.
\item[(c)] In particular, parallel turns bounding an annulus without
  singular vertex have the same shift for every $\ell$: a helix is a
  property of a whole cylinder of parallel turns, not of one of them.
\end{enumerate}
\end{proposition}

\begin{proof}
(a) Choose disjoint paths inside $D$ from $x_0 \in \bar\gamma$ to each
$c_i$. The curve $\bar\gamma$ is homotopic in $\Sigma^\circ$ to the product of
the small loops round the $c_i$ along these paths, in the order in which the
paths leave $x_0$. The holonomy being a homomorphism, $\hat h(\bar\gamma)$ is
the composition of the motions $x \mapsto R_i(x - p_i) + p_i$; its rotation is
$R^{\sum_i(4 - k(c_i))} = R^4 = I$ by Lemma~\ref{lem:gauss}, and expanding
the composition gives the sum. The shift is its component along $e_2$.
(b) Take $x_0$ on $\bar\gamma_1$ at the start of $\kappa$. The loop $\lambda
= \bar\gamma_1\, \kappa\, \bar\gamma_2^{-1} \kappa^{-1}$ is the border of
the disc obtained by cutting $A$ along $\kappa$. Developing along $\kappa$ is
a change of coordinates; conjugating the translation $\hat h(\bar\gamma_2)$
by it gives that translation read in the coordinates at $x_0$, where the walk
of $\gamma_2$ runs along $e_1$ by hypothesis. Hence $\hat h(\lambda)$ is the
translation by $\hat h(\bar\gamma_1) - \hat h(\bar\gamma_2)$, both read in
the same coordinates, whose component along $e_2$ is the difference of the
shifts, and part~(a) applied to this disc, whose singular
vertices have total index zero by Lemma~\ref{lem:gauss}, gives the sum.
(c) Without singular vertex the sum is empty.
\end{proof}

The simplest configuration between two parallel turns is a pair of
singular vertices of valences three and five, the defect that lets the number
of edges round a limb change. It is a dislocation of the grid.

\begin{proposition}[A dislocation between two turns]\label{prop:dislocation}
In the setting of Proposition~\ref{prop:burgers}(b), let the annulus $A$ hold
exactly two singular vertices, $c_3$ of valence three and $c_5$ of valence
five, and let $d = p_5 - p_3 = d_1 e_1 + d_2 e_2$ be their separation,
developed inside the cut disc. Then:
\begin{enumerate}
\item[(i)] $\shift_{\gamma_1}(\ell) - \shift_{\gamma_2}(\ell) =
  \varepsilon_1 d_1 + \varepsilon_2 d_2$ with signs $\varepsilon_1,
  \varepsilon_2 \in \{+1, -1\}$ that depend on the configuration but not on
  $\ell$, and $|B_A| = \sqrt2\,|d|$. The two turns can close together only
  if $c_5$ lies, as seen from $c_3$, on one particular diagonal of the field.
\item[(ii)] If $c_3$ and $c_5$ lie on one straight line of the field inside
  $A$, of length $\delta$ (in a T-mesh: they are joined by a chain of arcs
  along one trace, and $\delta$ is the sum of their lengths), then
  $|\shift_{\gamma_1}(\ell) - \shift_{\gamma_2}(\ell)| = \delta$. The two
  turns can both be rings only if $\delta = 0$, that is if the two
  singularities merge; as long as the bounds of separation keep
  $\delta \ge 1$, one of them at least is a helix.
\item[(iii)] For whole lengths, $\shift_{\gamma_1}(\ell) -
  \shift_{\gamma_2}(\ell) \equiv d_1 + d_2 \pmod 2$: a pair at an odd
  number of edges from one another, counted along the two directions of the
  field, never lies between two rings.
\end{enumerate}
\end{proposition}

\begin{proof}
With $R_3 = R$ and $R_5 = R^{-1}$, the sum of
Proposition~\ref{prop:burgers} for the two possible orders is
$B_A = R^{-1}(I - R)\,p_3 + (I - R^{-1})\,p_5 = (I - R^{-1})\,d$ or
$B_A = R\,(I - R^{-1})\,p_5 + (I - R)\,p_3 = (I - R)(-d)$. In coordinates
$(e_1, e_2)$, $R(a, b) = (-b, a)$, so
$(I - R^{-1})(d_1, d_2) = (d_1 - d_2,\ d_1 + d_2)$ and
$(I - R)(-d_1, -d_2) = (-d_1 - d_2,\ d_1 - d_2)$; reversing the
orientation of the border, which inverts the rotations, gives the two other
sign patterns. (We checked the
four cases symbolically.) In all of them both components of $B_A$ are of
the form $\pm d_1 \pm d_2$ and $|B_A|^2 = 2|d|^2$, which is (i); the
component along $e_2$ vanishes exactly on the line $\varepsilon_1 d_1 +
\varepsilon_2 d_2 = 0$, a diagonal. For (ii), choose $\kappa$ away from the
straight segment joining $c_3$ to $c_5$, which does not separate the two
borders of $A$, and develop $c_5$ from $c_3$ along it: $d$ is $\delta$ times
one of $\pm e_1$, $\pm e_2$, so $|\varepsilon_1 d_1 + \varepsilon_2 d_2| =
\delta$. For (iii), $d \in \Z^2$ and $\pm d_1 \pm d_2 \equiv d_1 + d_2$.
\end{proof}

Plain counting confirms (iii) independently. If both turns were rings, the
along component of $B_A$ would be the difference $n_1 - n_2$ of the numbers
of edges round them, of the parity of $d_1 + d_2$; but an annulus made of
quads has an even number of edges on its border, since $4F = 2E_{\text{inside}}
+ E_{\text{border}}$, so $n_1 + n_2$ is even. Part (ii) is the mechanism the
problem of feasibility in Section~\ref{sec:open} guessed: the conflict found on the male
body at a single arc of the left wrist would be explained if that arc joins
a vertex of valence three to one of valence five between two held rings. We
have not checked it.

For more singular vertices between two turns, the parity part generalizes
into a checkerboard rule.

\begin{proposition}[Checkerboard rule]\label{prop:checkerboard}
In the setting of Proposition~\ref{prop:burgers}(b), with whole lengths,
colour the points $p \in \Z^2$ of the cut disc black or white by the parity
of $p_1 + p_2$. Then $\shift_{\gamma_1}(\ell) - \shift_{\gamma_2}(\ell)$ has
the parity of the number of singular vertices of odd valence inside $A$
that develop onto white points. In particular, two rings bound an annulus
only if that number is even.
\end{proposition}

\begin{proof}
In the sum of Proposition~\ref{prop:burgers}, a vertex of even valence has
$R_i = I$ or $R_i = -I$, and $(I - R_i)p_i \in \{0, 2p_i\}$ is even. A vertex
of odd valence has $R_i = R^{\pm1}$, and both components of $(I - R^{\pm1})
(a, b) = (a \pm b,\ b \mp a)$ have the parity of $a + b$. Quarter-turns keep
the property ``both components have the parity $\sigma$'', and sums add the
$\sigma$. So both components of $B_A$ have the parity of the number of odd
vertices on white points. The count does not depend on the coordinates: a
quarter-turn keeps colours, and a whole translation $t$ changes the colour of
every point if $t_1 + t_2$ is odd, which keeps the parity of the count
because the number of odd vertices inside $A$ is even, their indices
$4 - k$ being odd and summing to zero with the even ones
(Lemma~\ref{lem:gauss}). Developing a vertex along another path inside the
cut disc moves it by a rotation $x \mapsto R_j(x - p_j) + p_j$ about another
singular vertex, which keeps its colour. We checked the rule on 13{,}000 random
configurations with sympy.
\end{proof}

\paragraph{How many independent rings.}
Rings that do not cross one another are disjoint simple closed curves of
$\Sigma^\circ$, none bounding a disc with fewer than two singular vertices
(Lemma~\ref{lem:gauss}). On a surface of genus zero with $n$ singular
vertices there are at most $n - 3$ such curves pairwise non-homotopic, the
number of curves of a decomposition of the $n$-punctured sphere into pairs of
pants; disjoint homotopic ones bound an annulus without singular vertex and
share their shift by Proposition~\ref{prop:burgers}(c). The rings of a
quad mesh that do not cross therefore fall into at most $n - 3$ cylinders,
the cylinders of the Jenkins--Strebel picture, and their closing conditions
have rank at most $n - 3$.

\begin{remark}[Whole periods are automatic; rings need zero]
For whole lengths, every shift is a whole number: the period of the integer
grid along a loop is always whole, which is what quantization guarantees and
what~\cite{kalberer2011stripe} enforces by rounding. Closing a ring is the
stronger requirement that this period be zero. The real shift of the
parametrization is the twist of the field's cylinder; it tells which whole
value rounding is likely to pick, not whether zero is reachable.
\end{remark}

\begin{remark}[Whole solutions exist at some size]
The balance and closing constraints are homogeneous and the bounds of
separation are lower bounds $\ell \ge b$ with $b \ge 0$. If they have a
real solution they have a rational one, a vertex of a rational polyhedron,
and multiplying it by the common denominator $N$ gives whole lengths
$N\ell \ge Nb \ge b$ that satisfy them all. Closing in whole numbers is thus
never impossible in itself; what can fail is closing \emph{at a given size}
(number of quads) or within a given deviation from the targets. The search for
integer solutions of Section~\ref{sec:open} should be read in that light.
\end{remark}

\begin{conjecture}[Conflicts are between neighbouring turns]\label{conj:local}
On a surface of genus zero, let $H$ be a set of turns each of which can be
closed on its own within the bounds of separation. If $H$ cannot be closed
together within those bounds, then $H$ contains two parallel turns whose
annulus has a Burgers vector with a component along $e_2$ that does not
vanish for any lengths balancing the patches and meeting the bounds.
\end{conjecture}

The converse holds by Proposition~\ref{prop:burgers}(b). The conjecture
says that the certificates of infeasibility given by linear programming
duality can always be taken of this local form. It is supported only by the
observation that the conflicts of Section~\ref{sec:exact} are local; it would
be refuted by an irreducible infeasible subset of the closing constraints
with more than two turns, which the exported system allows to search for.

\section{Open problems}\label{sec:open}

\begin{problem}[Exact formulation]
Minimize the deviation of the quantization subject to balance, validity and
$\shift_\gamma(\ell)=0$ for a set of turns. This is a minimum-deviation flow
with linear side constraints; how much of the exact machinery
of~\cite{heistermann2023bimdf} carries over is open. Section~\ref{sec:exact}
shows that the linear relaxation already closes almost all held turns.
\end{problem}

\paragraph{Exact whole lengths.}
Closing exactly, in whole numbers, the turns found closable together with
the bounds kept is feasible: HiGHS solves it in about two minutes for
three of the four ways tried (the fourth, with 87 turns, runs out of ten
minutes). The solutions leave 1 to 4 held turns shifted where the engine
leaves 13 to 44, and fewer of the turns not held as well (5 against 9, 5
against 11, 0 against 27), at a deviation from the real lengths 4\% to 62\%
higher. Yet the quad meshes they make spiral \emph{more}: 12.2\% of their
edges against 8.8\% for the first way, 3.7\% against 1.1\% for the second,
and the third breaks the structure. Most of the loops that spiral are thus
not turns our search finds: closing every turn it finds is not enough,
which makes the second open problem below the central one.

\paragraph{Closing by cuts.}
On the same body, walking every strip of the mesh the engine makes
(Proposition~\ref{prop:complete}) finds about 220 ways round per way of
closing, of which 12 to 24 come back shifted although the field itself
shifts them by less than half a row: the sampled walks of
Section~\ref{sec:closing} had found only 4 to 10 of those, and held 1 to 5.
A first selection of rings by that criterion (the field's shift under half
a row) proved too narrow: the twisted thigh, shifted 1.65 rows a turn by the
field, is a ring a modeller wants closed. Closing those rings by cuts
converges (12, then 14, then 6 shifted over three rounds) but leaves the
others free, and the mesh spirals more than with the engine (18\% of the
edges of that way against 1.1\%). Selecting rings by their lean instead, as
the engine does, gives 222 of the 230 ways round; 218 of them can be closed
together in real numbers with every bound of separation kept. In whole
numbers HiGHS found no solution within ten minutes. On this body, nearly
every loop could thus be a ring; what remains is integer.

\begin{problem}[Closing in whole numbers]
The solutions of the balance and shift equations in whole numbers form a
lattice. Can the integer search move within it, for instance from a
reduced basis (Hermite normal form, LLL) and a nearest-plane rounding of
the targets, instead of rounding first and repairing after? The
measurements above say this is where the loops are lost.
\end{problem}

\begin{problem}[Which turns]
The straight walks that return are infinite in number; we sample them, and
Section~\ref{sec:exact} shows the sample misses most of the loops that
spiral. Is there a finite set of turns, a basis in some homological sense
of the regular region of the T-mesh, whose closing implies that every loop
of the quad mesh closes? Can the loops of the quad mesh to be read off
the T-mesh and its lengths directly, so that every one is accounted for?
\end{problem}

\begin{problem}[Feasibility]
On a human torso, some held turns cannot be closed together: the projection
lets them go. Our hypothesis is that pairs of singularities between two
rings force a dislocation, and that closing all rings needs the poles to
face one another, which the separation of singularities forbids at the
scale of the T-mesh; Proposition~\ref{prop:dislocation} proves this for a
single pair on a trace. A capped cylinder shows the same conflict: its
layout of six patches exists, but only with arcs our quantizer will not
bring to zero. On the male body the conflict comes down to one or a few
arcs (Section~\ref{sec:exact}). When is a set of turns simultaneously
closable, in terms of the singularity configuration?
Conjecture~\ref{conj:local} proposes that the obstructions are always
between two parallel turns.
\end{problem}

\begin{problem}[Finite trajectories]
In the language of~\cite{lei2020quartic}: given a singularity configuration
on a surface of genus zero, when does a quartic differential with all
trajectories finite exist, and how close to a given cross field can it be
chosen? Our turns are its discrete, integer version.
\end{problem}

\paragraph{Towards proofs.}
Proposition~\ref{prop:linear} is elementary but its statement over a
T-mesh with T-junctions is easy to get wrong. A formal proof in a proof
assistant, over an abstract complex of rectangles with integer side
lengths, would settle it and could carry the sufficiency question of the
second problem for special classes of layouts.

\paragraph{Reproducibility.}
The method is implemented in openclare, Rust, branches
\texttt{close-quantized} (closing, measure) and \texttt{quadsphere}
(confidence floor, merged traces); the bench is driven by
\texttt{bench/run.py} and the marks by \texttt{scripts/scores.py}.

\bibliographystyle{plainnat}
\bibliography{refs}

\end{document}
